% answers_short.tex —— 2018 年·江苏卷 · 简短答案速查
% 由 tools/make_short_answers.py 自 answers.tex 的 \gksolhead/\gkanswer 提取，
% 并与 2018/2018_数学真题_江苏_答案要点.md 逐题比对后生成，勿手改（改源 MD 后重跑）。

\begin{gknotice}
\textbf{使用说明}：本节只给最终答案便于快速核对，完整推导见解析版。答案为编写者独立推导，\textbf{非官方标准答案}。
\end{gknotice}

\gkpart{填空题}
\begin{itemize}[leftmargin=2.6em]
  \item[\textbf{1.}] $\{1, 8\}$
  \item[\textbf{2.}] $2$
  \item[\textbf{3.}] $90$
  \item[\textbf{4.}] $8$
  \item[\textbf{5.}] $[2, +\infty)$
  \item[\textbf{6.}] $\dfrac{3}{10}$
  \item[\textbf{7.}] $-\dfrac{\pi}{6}$
  \item[\textbf{8.}] $2$
  \item[\textbf{9.}] $0$
  \item[\textbf{10.}] $\dfrac{4}{3}$
  \item[\textbf{11.}] $-3$
  \item[\textbf{12.}] $3$
  \item[\textbf{13.}] $9$
  \item[\textbf{14.}] $27$
\end{itemize}

\gkpart{解答题}
\begin{itemize}[leftmargin=2.6em]
  \item[\textbf{15.}] (1) 见解析版；(2) 见解析版
  \item[\textbf{16.}] (1) $-\dfrac{7}{25}$；(2) $-\dfrac{2}{11}$
  \item[\textbf{17.}] (1) 矩形 $ABCD$ 的面积为 $800(4\sin\theta\cos\theta+\cos\theta)$ 平方米，$\triangle CDP$ 的面积为 $1600(\cos\theta-\sin\theta\cos\theta)$ 平方米，$\sin\theta$ 的取值范围是 $\left[\dfrac{1}{4}, 1\right)$；(2) 当 $\theta = \dfrac{\pi}{6}$ 时，能使甲、乙两种蔬菜的年总产值最大
  \item[\textbf{18.}] (1) 椭圆 $C: \dfrac{x^2}{4}+y^2=1$，圆 $O: x^2+y^2=3$；(2) ① 点 $P$ 的坐标为 $(\sqrt{2}, 1)$；② 直线 $l$ 的方程为 $\sqrt{5}x+y-3\sqrt{2}=0$（或 $\sqrt{10}x+\sqrt{2}y-6=0$）
  \item[\textbf{19.}] (1) 见解析版；(2) $a = \dfrac{1}{2\mathrm{e}}$；(3) 存在
  \item[\textbf{20.}] (1) $\left[\dfrac{7}{3}, \dfrac{5}{2}\right]$；(2) $d$ 的取值范围为 $\left[\dfrac{b_1(q^m-2)}{m}, \dfrac{b_1 q^m}{m}\right]$（存在性证明见解析版）
\end{itemize}

\gkpart{附加题}
\begin{itemize}[leftmargin=2.6em]
  \item[\textbf{21.}] 【A】$BC=2$；【B】$A^{-1}=\begin{bmatrix}2 & -3 \\ -1 & 2\end{bmatrix}$，点 $P$ 坐标为 $(3, -1)$；【C】$2\sqrt{3}$；【D】$4$。
  \item[\textbf{22.}] (1) $\dfrac{3\sqrt{10}}{20}$；(2) $\dfrac{\sqrt{5}}{5}$
  \item[\textbf{23.}] (1) $f_3(2)=2$，$f_4(2)=5$；(2) $f_n(2)=\dfrac{n^2-n-2}{2}$ ($n\geqslant 5$)
\end{itemize}
